\section{Concluding remarks}\label{sec:conclusion}

The three-dimensional argument consists of a geometric reduction followed by a polynomial calculation. The generic-fiber construction uses the trace and principal-minor identities to obtain a plane equation and a quadratic equation. Homogeneous triangularization then restricts the directions at infinity. Once a dependent pair is available, the normal form and its tame factorization follow from the derivation argument of Proposition~\ref{prop:pair-normal}. The descent uses the one-dimensional space of quadratic relations, so the final conjugating maps are defined over the original field.

The canonical family also gives several explicit consequences. If $d=\deg P\geq1$, then
\[
 \deg\Ncal_P=4d,\qquad \rank J\Ncal_P=2.
\]
The minor in rows $1,2$ and columns $2,3$ of~\eqref{eq:canonical-J} is $P'(y+x^2)\neq0$. The $(1,3)$ entry of $(J\Ncal_P)^2$ is the same nonzero polynomial. Thus the nilpotency index is exactly three. The classification therefore implies that every normalized map in dimension three with nilpotent Jacobian and linearly independent components has degree divisible by four. This is a consequence of the proposed general theorem, whereas the displayed degree, rank, and nilpotency-index calculations for the canonical family are direct.

For the block extensions, the distinction between coefficients and polynomial variables is essential. The last three components must first be normalized over the fraction field of the base ring. The resulting linear normal form need not arise from a change of coordinates over the polynomial ring. Stable tameness follows from a separate argument: polynomial inversion descends through the formal inverse at the origin, and residue-field tameness permits the application of the regular-ring criterion. This argument does not require descent of the rational conjugating matrices.

The recent work of Casta{\~n}eda, Honorato, and Valenzuela-Henr{\'i}quez~\cite[Theorems~1.2, 5.1, and~8.1]{CHV2026} connects Keller maps with the weak Markus--Yamabe injectivity problem. Starting from a real Keller map $F=\id+H$ with $H(0)=JH(0)=0$ and component degrees $d_i$, their chain realization constructs a polynomial vector field $X_{F,\beta}$ satisfying
\[
 N=\sum_{i=1}^{n}\max\{d_i,2\}-n,
 \qquad
 \operatorname{Spec}\bigl(JX_{F,\beta}(Z)\bigr)=\{-1\}
 \quad (Z\in\mathbb R^N).
\]
Its equilibria correspond polynomially to $F^{-1}(\beta)$. Component degrees $(4,6,7)$ yield a degree-seven vector field in dimension $14$ with three rational equilibria. A separate construction yields a degree-three vector field in dimension $18$. Adjoining negative identity vector fields extends the counterexamples to every dimension at least $14$.

For these fields, $R=\id+X$ has nilpotent Jacobian, whereas $\id-R=-X$ is noninjective. Thus nilpotence alone does not imply invertibility in arbitrary dimension. The constructions do not impose the variable dependence required by the block theorems of Section~\ref{sec:blocks}.

In dimension three, Theorem~\ref{thm1} gives injectivity for the constant-spectrum class $X=-\id+H$ with $JH$ nilpotent. This does not settle the general Hurwitz case, where the eigenvalues need only have negative real parts. The cited work leaves dimensions $3$ through $13$ unresolved.

Several questions lie beyond the scope of these arguments. In dimension four, can the nilpotence identities control the degree of the embedded generic curve when the Jacobian rank is three? Under which additional hypotheses can two independent constant linear combinations of the components be made algebraically dependent? For the block extensions, can one bound the required stabilization in terms of the dimension or degrees? Answering these questions requires additional arguments; the three-dimensional plane-conic calculation does not extend formally to a general higher-dimensional map.
